\section{Introduction}
\label{sec:introduction}

Long-context language models place acute pressure on the key--value (KV) cache, which grows linearly with sequence length \cite{kwon-sosp2023, pope-mlsys2023}.
Two complementary responses---\emph{attention head pruning} \cite{voita-acl2019, michel-neurips2019} and \emph{KV cache eviction} \cite{snapkv-neurips2024, adakv-neurips2025}---are increasingly deployed simultaneously as \emph{compound sparsity}, but the resulting quality degradation is often super-additive \cite{spark-aaai2026, slim-icml2025}. Recent evaluations show that reasoning capabilities degrade disproportionately under compound compression \cite{unicomp-2026}, and that per-dimension safe zones intersect to form a smaller joint safe region \cite{phase-transition-compression-2026}.

Prior work treats compound interaction magnitude as roughly constant across generation \cite{spark-aaai2026, unicomp-2026}, despite evidence that pruning sensitivity depends on representational phase \cite{decision-repr-transition-2026} and that no single sparse attention strategy dominates across generation stages \cite{sparse-frontier-2025}.
This assumption is problematic for reasoning models, whose chain-of-thought generations \cite{wei-neurips2022-cot} alternate between broadly distributed attention (exploratory reasoning) and sharply concentrated attention (structured computation) \cite{thinkv-iclr2026}.
Phase-aware methods like ThinKV \cite{thinkv-iclr2026} apply different KV strategies per phase but without considering compound interaction, while importance-based eviction itself faces challenges from score collapse \cite{hies-2025}, greedy bias \cite{greedy-bias-kvcache-2026}, and temporal dynamics \cite{lethe-aaai2026}. These methods each optimize a single pruning axis; we characterize the interaction when multiple axes are composed.
If compound severity varies across these phases, a uniform sparsity budget may be safe in one phase yet destructive in another.

We find that compound sparsity severity is an emergent property of the three-way interaction among eviction strategy, reasoning phase, and sparsity level.
Under importance-based KV cache eviction, the Interaction Index---a measure of super-additive degradation---concentrates in low-sparsity reasoning phases and is largely absent in high-sparsity phases.
On DeepSeek-R1-Distill-Qwen-7B (MATH-500), the paired median difference rises from $0.39$ at $s{=}0.15$ to $8.48$ at $s{=}0.30$ (Spearman $\rho{=}1.0$; all Bonferroni-corrected $p < 0.005$).
The pattern replicates on Qwen2.5-7B-Instruct ($\Delta$: $0.07 \to 3.16$; all $p < 0.005$), with cross-architecture support from DeepSeek-R1-Distill-Llama-8B (LlamaForCausalLM, $n{=}250$), where the importance phase gradient replicates with large effect sizes ($d_z{=}0.81$--$1.15$; all $p < 10^{-29}$) despite reduced strategy specificity ($2.0$--$2.6{\times}$ vs.\ $15$--$17{\times}$ on Qwen2-family models).
Random eviction develops a detectable phase gradient only above $s \approx 0.20$, at $2.7$--$17{\times}$ smaller magnitude than importance-based (where both strategies reach significance at $s \geq 0.25$), supporting that the dominant gradient is strategy-specific.
A scoring-independent control (recency eviction, which removes the oldest positions without attention scores) produces no phase gradient ($\Delta{=}-0.003$, CI spans zero), ruling out segmentation artifact and confirming that the interaction is specifically driven by attention-score-based selection (\S\ref{sec:scoring-control}).
Cross-benchmark validation on GSM8K reveals consistent phase gradient direction: DeepSeek replicates the monotonic dose-response ($\Delta$: $0.24 \to 4.68$; all $p \leq 0.0003$), while Qwen shows directional consistency at $s \geq 0.25$ but is inconclusive due to generation truncation (\S\ref{sec:cross-benchmark}). Strategy specificity attenuates on GSM8K relative to MATH-500, indicating that phase modulation generalizes but its magnitude is task-dependent.

Per-phase Gini analysis identifies \emph{importance discrimination collapse}---the degeneration of head importance scores toward uniformity under compound stress---as a correlational diagnostic of this phase--strategy divergence. Both strategies exhibit consistent Gini drops across all tested conditions, yet severity diverges by $2.7$--$17{\times}$, implicating the structural specificity of importance-based pruning as the key differentiator.

Beyond individual findings, we contribute a \emph{phase-gated analysis framework} that decomposes compound severity by reasoning phase and eviction strategy, transforming a scalar interaction metric into a structured diagnostic. This framework---combining the Interaction Index with attention-sparsity-based phase segmentation---reveals structure in compound degradation that aggregate metrics miss, and provides practitioners with a principled basis for phase-aware sparsity allocation.

Our contributions:
\begin{enumerate}[nosep,leftmargin=*]
    \item We characterize a \textbf{phase-modulated compound severity gradient}: importance-based eviction exhibits a monotonic dose-response ($\rho{=}1.0$; $\Delta$: $0.39 \to 8.48$) with a phase gradient $2.7$--$17{\times}$ larger than random eviction where both strategies are significant ($s \geq 0.25$) on MATH-500; the gradient direction is consistent on GSM8K, while strategy specificity attenuates (\S\ref{sec:cross-benchmark}).
    \item We replicate the pattern \textbf{across training paradigms and architectures}: within-architecture replication on Qwen2.5-7B-Instruct (same Qwen2 backbone, instruction-tuned; $\Delta$: $0.07 \to 3.16$; $\rho{=}1.0$) and cross-architecture validation on DeepSeek-R1-Distill-Llama-8B (LlamaForCausalLM; $n{=}250$; importance phase gradient $d_z{=}0.81$--$1.15$, all $p < 10^{-29}$; strategy specificity reduced to $2.0$--$2.6{\times}$, architecture-dependent; \S\ref{sec:cross-arch}).
    \item We provide a \textbf{scoring-independent control}: recency eviction (position-only, no attention scores) under identical compound sparsity produces no phase gradient ($\Delta \approx 0$, CI spans zero), ruling out segmentation artifact and establishing that phase-dependent severity is specifically driven by attention-score-based selection (\S\ref{sec:scoring-control}).
    \item We identify \textbf{importance discrimination collapse} via per-phase Gini analysis as a correlational diagnostic: consistent collapse across models and strategies, with the structural specificity of importance-based pruning corresponding to the severity divergence (\S\ref{sec:analysis}).
    \item We validate the phase gradient on \textbf{GSM8K} for both models, finding that phase modulation is benchmark-robust while strategy specificity is task-dependent (\S\ref{sec:cross-benchmark}).
\end{enumerate}
